\documentclass[12pt,a4paper]{report}

\usepackage[utf8]{inputenc}
\usepackage[T1]{fontenc}
\usepackage[french]{babel}
\usepackage{geometry}
\usepackage{longtable}
\usepackage{array}
\usepackage{xcolor}
\usepackage{hyperref}
\usepackage{enumitem}
\usepackage{listings}
\usepackage{titlesec}

\geometry{margin=2.3cm}
\hypersetup{
    colorlinks=true,
    linkcolor=blue,
    urlcolor=blue
}

\definecolor{codegray}{RGB}{245,245,245}
\definecolor{darkblue}{RGB}{25,55,95}

\lstset{
    basicstyle=\ttfamily\small,
    backgroundcolor=\color{codegray},
    breaklines=true,
    frame=single,
    columns=fullflexible
}

\titleformat{\chapter}{\normalfont\LARGE\bfseries\color{darkblue}}{\thechapter}{1em}{}
\titleformat{\section}{\normalfont\Large\bfseries\color{darkblue}}{\thesection}{1em}{}
\titleformat{\subsection}{\normalfont\large\bfseries}{\thesubsection}{1em}{}

\title{
    \textbf{Documentation technique du backend SmartRecruit}\\
    \large Système de classement de CV par fiche de poste
}
\author{Projet PFE -- Backend FastAPI, RAG, Qwen et PostgreSQL}
\date{\today}

\begin{document}

\maketitle
\tableofcontents
\newpage

\chapter{Présentation générale du backend}

Le backend de \textbf{SmartRecruit} est une application FastAPI qui analyse une fiche de poste et plusieurs CV afin de produire un classement explicable des candidats. Le système ne se limite pas à une comparaison simple de mots-clés. Il combine plusieurs niveaux de traitement :

\begin{itemize}
    \item extraction du texte depuis des documents PDF, DOCX, TXT ou MD ;
    \item segmentation du texte en sections utiles ;
    \item extraction structurée de la fiche de poste et des CV ;
    \item normalisation des compétences, langues, diplômes et intitulés ;
    \item indexation sémantique des sections de CV ;
    \item recherche RAG avec embeddings Qwen et stockage vectoriel PostgreSQL ;
    \item calcul de scores par catégories ;
    \item classement final avec preuves, forces, faiblesses et gestion des égalités.
\end{itemize}

L'objectif principal est de comparer chaque candidat à une fiche de poste précise. La fiche de poste sert de référence d'évaluation. Les CV représentent les éléments à analyser, vectoriser, comparer et classer.

\section{Entrées et sorties}

\subsection{Entrées}

Le backend reçoit principalement :

\begin{enumerate}
    \item une fiche de poste, appelée aussi fiche de test ou fiche de critères ;
    \item un ou plusieurs CV candidats ;
    \item un paramètre optionnel \texttt{top\_k}, qui indique combien de preuves sémantiques doivent être récupérées par candidat.
\end{enumerate}

\subsection{Sorties}

La réponse contient :

\begin{itemize}
    \item la fiche de poste structurée ;
    \item le nombre total de candidats analysés ;
    \item le classement final ;
    \item pour chaque candidat : score final, scores par catégorie, preuves, forces, faiblesses, CV structuré ;
    \item les erreurs éventuelles par fichier.
\end{itemize}

\chapter{Architecture logique}

\section{Vue globale}

Le déroulement global peut être résumé ainsi :

\begin{lstlisting}
Utilisateur
   |
   |-- Upload fiche de poste
   |-- Upload un ou plusieurs CV
   v
FastAPI
   |
   |-- Extraction du texte
   |-- Segmentation en sections
   |-- Extraction structuree CV / fiche de poste
   |-- Normalisation
   |-- Indexation vectorielle des CV
   |-- Recherche semantique RAG
   |-- Matching et scoring
   |-- Classement final
   v
Reponse JSON explicable
\end{lstlisting}

\section{Pipeline détaillé}

\subsection{Étape 1 : réception des fichiers}

L'utilisateur envoie une fiche de poste et plusieurs CV via l'endpoint principal :

\begin{lstlisting}
POST /api/ranking/analyze
\end{lstlisting}

Les fichiers sont transmis en \texttt{multipart/form-data}. Le backend vérifie leur format et les transmet au pipeline de classement.

\subsection{Étape 2 : extraction documentaire}

Le service \texttt{DoclingParser} lit le contenu du document selon son format :

\begin{itemize}
    \item PDF avec \texttt{PyMuPDF} ;
    \item DOCX avec \texttt{python-docx} ;
    \item TXT ou MD avec lecture texte classique.
\end{itemize}

Le résultat est représenté par un objet \texttt{DocumentText}, contenant :

\begin{itemize}
    \item le nom du fichier ;
    \item le type de document ;
    \item le texte brut ;
    \item le nombre de caractères ;
    \item les sections détectées.
\end{itemize}

\subsection{Étape 3 : segmentation en sections}

Le backend découpe le texte en sections comme :

\begin{itemize}
    \item compétences ;
    \item expériences ;
    \item formation ;
    \item langues ;
    \item projets ;
    \item responsabilités ;
    \item exigences.
\end{itemize}

Cette étape est importante pour éviter de mélanger formation, langues, projets et expériences professionnelles.

\subsection{Étape 4 : extraction structurée}

Deux extracteurs principaux sont utilisés :

\begin{itemize}
    \item \texttt{JobExtractor} pour structurer la fiche de poste ;
    \item \texttt{CVExtractor} pour structurer chaque CV.
\end{itemize}

Le backend utilise Qwen via un endpoint compatible OpenAI. Les extracteurs appellent le modèle configuré pour produire un JSON structuré, puis le backend valide et normalise ce JSON avant le scoring.

\subsection{Étape 5 : normalisation}

Les données extraites sont normalisées pour éviter que des variantes d'écriture faussent le matching. Par exemple :

\begin{itemize}
    \item \texttt{PowerBI} devient \texttt{power bi} ;
    \item \texttt{dashbord} devient \texttt{dashboard} ;
    \item \texttt{Postgres} devient \texttt{postgresql} ;
    \item les niveaux de langue sont ramenés à une échelle comparable ;
    \item les diplômes sont convertis en niveaux hiérarchiques.
\end{itemize}

\subsection{Étape 6 : indexation vectorielle}

Les sections des CV sont découpées en chunks. Chaque chunk est transformé en vecteur par le client d'embeddings Qwen. Ces vecteurs sont stockés dans PostgreSQL avec SQLAlchemy.

\subsection{Étape 7 : recherche RAG}

Le backend construit une requête à partir de la fiche de poste :

\begin{itemize}
    \item intitulé du poste ;
    \item compétences obligatoires ;
    \item compétences souhaitées ;
    \item responsabilités nettoyées.
\end{itemize}

Cette requête est vectorisée, puis utilisée pour récupérer les passages les plus pertinents du CV dans PostgreSQL par similarité cosinus entre vecteurs Qwen.

\subsection{Étape 8 : matching et scoring}

Le scoring est divisé en catégories :

\begin{itemize}
    \item compétences techniques ;
    \item expérience pertinente ;
    \item responsabilités ;
    \item formation ;
    \item langues ;
    \item soft skills.
\end{itemize}

Les catégories absentes dans la fiche de poste ne donnent pas automatiquement 100/100. Elles sont ignorées et leurs poids sont redistribués sur les catégories réellement applicables.

\subsection{Étape 9 : classement}

Les candidats sont triés par score final décroissant. Si deux candidats ont exactement le même score, le système indique une égalité avec \texttt{ex aequo}.

\chapter{Rôle des dossiers du backend}

\begin{longtable}{|p{4cm}|p{10cm}|}
\hline
\textbf{Dossier} & \textbf{Rôle} \\
\hline
\endhead
\texttt{backend/app} & Contient le code principal de l'application FastAPI. \\
\hline
\texttt{backend/app/api} & Regroupe les routes HTTP exposées par l'API. \\
\hline
\texttt{backend/app/api/routes} & Contient les endpoints concrets : santé, documents, classement. \\
\hline
\texttt{backend/app/core} & Contient les éléments transversaux : constantes, exceptions, configuration du logging. \\
\hline
\texttt{backend/app/data} & Contient les fichiers JSON de référence : alias de compétences, niveaux de diplômes, niveaux de langues, poids du scoring. \\
\hline
\texttt{backend/app/database} & Définit la base SQLAlchemy et la table PostgreSQL utilisée pour stocker les chunks vectorisés. \\
\hline
\texttt{backend/app/infrastructure} & Contient les clients techniques vers Qwen et PostgreSQL : LLM, embeddings et stockage vectoriel SQLAlchemy. \\
\hline
\texttt{backend/app/schemas} & Définit les modèles Pydantic utilisés dans l'application : CV, fiche de poste, matching, ranking, document. \\
\hline
\texttt{backend/app/services} & Contient la logique métier principale : extraction, normalisation, RAG, matching, scoring et orchestration. \\
\hline
\texttt{backend/scripts} & Contient les scripts d'aide pour lancer ou préparer l'environnement : GPU, vLLM/Qwen, bases et backend. \\
\hline
\texttt{backend/tests} & Contient les tests unitaires et d'intégration du backend. \\
\hline
\texttt{backend/uploads} & Dossier prévu pour recevoir temporairement les fichiers uploadés. \\
\hline
\texttt{backend/storage} & Dossier prévu pour stocker documents et résultats. \\
\hline
\end{longtable}

\chapter{Rôle des fichiers principaux}

\section{Racine du backend}

\begin{longtable}{|p{5cm}|p{9cm}|}
\hline
\textbf{Fichier} & \textbf{Rôle} \\
\hline
\endhead
\texttt{backend/requirements.txt} & Liste les dépendances Python nécessaires : FastAPI, Uvicorn, PyMuPDF, SQLAlchemy, psycopg2, pytest, etc. \\
\hline
\texttt{backend/docker-compose.yml} & Prépare le service PostgreSQL local. \\
\hline
\texttt{backend/app/main.py} & Point d'entrée FastAPI. Il crée l'application, configure CORS et inclut les routes. \\
\hline
\texttt{backend/app/config.py} & Centralise la configuration : chemins, URL Qwen, modèles Qwen, PostgreSQL et limite d'upload. \\
\hline
\texttt{backend/app/dependencies.py} & Assemble les dépendances principales : parser, client LLM Qwen, client embeddings Qwen, stockage vectoriel PostgreSQL et pipeline de classement. \\
\hline
\end{longtable}

\section{Routes API}

\begin{longtable}{|p{5cm}|p{9cm}|}
\hline
\textbf{Fichier} & \textbf{Rôle} \\
\hline
\endhead
\texttt{api/routes/health.py} & Expose une route simple de santé pour vérifier que le backend fonctionne. \\
\hline
\texttt{api/routes/documents.py} & Prépare les routes liées aux documents. \\
\hline
\texttt{api/routes/ranking.py} & Expose l'endpoint principal \texttt{/api/ranking/analyze}. Il reçoit la fiche de poste et les CV, puis lance le pipeline de classement. \\
\hline
\end{longtable}

\section{Schemas Pydantic}

\begin{longtable}{|p{5cm}|p{9cm}|}
\hline
\textbf{Fichier} & \textbf{Rôle} \\
\hline
\endhead
\texttt{schemas/document.py} & Définit \texttt{DocumentText}, structure représentant le texte extrait d'un document. \\
\hline
\texttt{schemas/cv.py} & Définit la structure d'un CV : compétences, expériences, formation, langues, projets, certifications. \\
\hline
\texttt{schemas/job.py} & Définit la fiche de poste structurée : compétences requises, expérience, formation, langues, responsabilités. \\
\hline
\texttt{schemas/experience.py} & Définit les objets liés aux durées d'expérience et aux périodes calculées. \\
\hline
\texttt{schemas/matching.py} & Définit les résultats de matching : scores par catégorie, preuves, score final candidat. \\
\hline
\texttt{schemas/ranking.py} & Définit la réponse finale de classement et la gestion des rangs, y compris les égalités. \\
\hline
\end{longtable}

\section{Services documents}

\begin{longtable}{|p{5cm}|p{9cm}|}
\hline
\textbf{Fichier} & \textbf{Rôle} \\
\hline
\endhead
\texttt{documents/docling\_parser.py} & Extrait le texte des PDF, DOCX, TXT ou MD et retourne un objet \texttt{DocumentText}. \\
\hline
\texttt{documents/section\_segmenter.py} & Détecte les sections d'un document et conserve leur contenu avec les retours ligne. Cela aide à distinguer expérience, formation, langues et projets. \\
\hline
\end{longtable}

\section{Services extraction}

\begin{longtable}{|p{5cm}|p{9cm}|}
\hline
\textbf{Fichier} & \textbf{Rôle} \\
\hline
\endhead
\texttt{extraction/cv\_extractor.py} & Extrait un CV structuré. Il filtre les faux intitulés d'expérience, récupère certaines entreprises connues, calcule les formations et évite de mélanger formation/projet/expérience. \\
\hline
\texttt{extraction/job\_extractor.py} & Extrait la fiche de poste structurée. Il nettoie les responsabilités pour ne garder que les vraies missions et sépare les compétences des responsabilités. \\
\hline
\texttt{extraction/prompts.py} & Contient les prompts utilisés pour demander au LLM une sortie JSON structurée. \\
\hline
\texttt{extraction/output\_validator.py} & Valide les sorties du LLM avec les modèles Pydantic attendus. \\
\hline
\end{longtable}

\section{Services normalisation}

\begin{longtable}{|p{5cm}|p{9cm}|}
\hline
\textbf{Fichier} & \textbf{Rôle} \\
\hline
\endhead
\texttt{normalization/text\_normalizer.py} & Nettoie le texte : minuscules, suppression des accents, espaces, ponctuation. \\
\hline
\texttt{normalization/skill\_normalizer.py} & Normalise les compétences avec un dictionnaire d'alias. \\
\hline
\texttt{normalization/job\_title\_normalizer.py} & Normalise les intitulés de poste pour comparer les expériences au poste demandé. \\
\hline
\texttt{normalization/education\_normalizer.py} & Normalise les diplômes et calcule leur rang. \\
\hline
\texttt{normalization/language\_normalizer.py} & Normalise les langues et les niveaux linguistiques. \\
\hline
\texttt{normalization/date\_normalizer.py} & Interprète les dates et les mois, y compris les abréviations françaises comme \texttt{janv}, \texttt{juil}, \texttt{dec}. \\
\hline
\end{longtable}

\section{Services expérience}

\begin{longtable}{|p{5cm}|p{9cm}|}
\hline
\textbf{Fichier} & \textbf{Rôle} \\
\hline
\endhead
\texttt{experience/duration\_calculator.py} & Calcule la durée d'une expérience à partir de dates ou d'une durée déclarée. \\
\hline
\texttt{experience/overlap\_manager.py} & Fusionne les périodes qui se chevauchent pour éviter de compter deux fois la même période. \\
\hline
\texttt{experience/relevance\_calculator.py} & Évalue si une expérience est pertinente pour la fiche de poste. Une expérience longue mais non pertinente ne doit pas donner un score artificiel. \\
\hline
\end{longtable}

\section{Services retrieval et RAG}

\begin{longtable}{|p{5cm}|p{9cm}|}
\hline
\textbf{Fichier} & \textbf{Rôle} \\
\hline
\endhead
\texttt{retrieval/chunk\_builder.py} & Découpe les sections en chunks utilisables pour l'indexation vectorielle. \\
\hline
\texttt{retrieval/section\_indexer.py} & Transforme les chunks en embeddings puis les insère dans le stockage vectoriel PostgreSQL. \\
\hline
\texttt{retrieval/semantic\_retriever.py} & Recherche les passages pertinents dans le stockage vectoriel PostgreSQL à partir de la requête issue de la fiche de poste. \\
\hline
\end{longtable}

\section{Services matching}

\begin{longtable}{|p{5cm}|p{9cm}|}
\hline
\textbf{Fichier} & \textbf{Rôle} \\
\hline
\endhead
\texttt{matching/skill\_matcher.py} & Compare les compétences du CV avec les compétences obligatoires et souhaitées. Il sépare \texttt{missing\_mandatory} et \texttt{missing\_preferred}. \\
\hline
\texttt{matching/experience\_matcher.py} & Compare l'expérience pertinente du candidat avec l'expérience demandée. \\
\hline
\texttt{matching/responsibility\_matcher.py} & Évalue les responsabilités. Il distingue preuve complète, preuve partielle et absence de preuve. \\
\hline
\texttt{matching/education\_matcher.py} & Compare le niveau de formation du candidat avec le niveau demandé. \\
\hline
\texttt{matching/language\_matcher.py} & Compare les langues et niveaux demandés avec ceux du candidat. \\
\hline
\texttt{matching/certification\_matcher.py} & Prépare la comparaison des certifications. \\
\hline
\end{longtable}

\section{Services scoring et ranking}

\begin{longtable}{|p{5cm}|p{9cm}|}
\hline
\textbf{Fichier} & \textbf{Rôle} \\
\hline
\endhead
\texttt{scoring/weights.py} & Charge les poids des catégories depuis \texttt{scoring\_weights.json}. \\
\hline
\texttt{scoring/scoring\_engine.py} & Calcule le score final. Il ignore les catégories absentes et redistribue les poids entre les critères réellement présents. \\
\hline
\texttt{scoring/explanation\_builder.py} & Produit les forces et faiblesses à partir des scores par catégorie. \\
\hline
\texttt{ranking/ranking\_engine.py} & Trie les candidats, attribue les rangs et marque les égalités avec \texttt{ex aequo}. \\
\hline
\end{longtable}

\section{Services orchestration}

\begin{longtable}{|p{5cm}|p{9cm}|}
\hline
\textbf{Fichier} & \textbf{Rôle} \\
\hline
\endhead
\texttt{orchestration/analyze\_job\_pipeline.py} & Exécute l'analyse complète d'une fiche de poste. \\
\hline
\texttt{orchestration/analyze\_cv\_pipeline.py} & Exécute l'analyse complète d'un CV. \\
\hline
\texttt{orchestration/batch\_ranking\_pipeline.py} & Coordonne tout le processus : analyse fiche, analyse CV, indexation, retrieval, scoring et classement. \\
\hline
\end{longtable}

\section{Infrastructure}

\begin{longtable}{|p{5cm}|p{9cm}|}
\hline
\textbf{Fichier} & \textbf{Rôle} \\
\hline
\endhead
\texttt{infrastructure/qwen\_llm.py} & Client LLM Qwen via endpoint compatible OpenAI. Utilisé pour générer du JSON structuré. \\
\hline
\texttt{infrastructure/qwen\_embeddings.py} & Client d'embeddings Qwen. Transforme les chunks et requêtes en vecteurs. \\
\hline
\texttt{infrastructure/postgres\_vector\_store.py} & Stocke les chunks vectorisés dans PostgreSQL avec SQLAlchemy et récupère les passages par similarité cosinus. \\
\hline
\end{longtable}

\section{Données JSON}

\begin{longtable}{|p{5cm}|p{9cm}|}
\hline
\textbf{Fichier} & \textbf{Rôle} \\
\hline
\endhead
\texttt{data/skill\_aliases.json} & Contient les alias de compétences : \texttt{PowerBI}, \texttt{dashbord}, \texttt{Postgres}, etc. \\
\hline
\texttt{data/job\_title\_aliases.json} & Contient les variantes d'intitulés de poste. \\
\hline
\texttt{data/education\_levels.json} & Définit la hiérarchie des diplômes. \\
\hline
\texttt{data/language\_levels.json} & Définit la hiérarchie des niveaux de langues. \\
\hline
\texttt{data/scoring\_weights.json} & Définit les poids de base des catégories de scoring. \\
\hline
\end{longtable}

\section{Tests}

\begin{longtable}{|p{5cm}|p{9cm}|}
\hline
\textbf{Fichier} & \textbf{Rôle} \\
\hline
\endhead
\texttt{tests/conftest.py} & Ajoute le backend au chemin Python pour exécuter les tests. \\
\hline
\texttt{tests/test\_health.py} & Vérifie que l'API répond correctement. \\
\hline
\texttt{tests/integration/test\_api\_ranking.py} & Teste l'endpoint de classement avec upload de fichiers. \\
\hline
\texttt{tests/integration/test\_ranking\_pipeline.py} & Teste le pipeline complet de classement. \\
\hline
\texttt{tests/unit/test\_experience.py} & Vérifie les calculs de durée et les dates. \\
\hline
\texttt{tests/unit/test\_extraction.py} & Vérifie l'extraction des fiches, CV, formations, entreprises et expériences. \\
\hline
\texttt{tests/unit/test\_matching.py} & Vérifie les règles de matching. \\
\hline
\texttt{tests/unit/test\_normalization.py} & Vérifie les fonctions de normalisation. \\
\hline
\texttt{tests/unit/test\_ranking.py} & Vérifie la gestion des rangs et égalités. \\
\hline
\texttt{tests/unit/test\_scoring.py} & Vérifie le scoring, les responsabilités, les preuves et les catégories manquantes. \\
\hline
\end{longtable}

\chapter{Explication du scoring}

\section{Compétences techniques}

Les compétences sont séparées en deux groupes :

\begin{itemize}
    \item compétences obligatoires ;
    \item compétences souhaitées.
\end{itemize}

Le score donne plus d'importance aux compétences obligatoires. Les compétences manquantes sont aussi séparées :

\begin{lstlisting}
missing_mandatory = ["power bi", "excel", "snowflake"]
missing_preferred = ["foundry", "spm"]
\end{lstlisting}

Ainsi, l'absence d'une compétence obligatoire coûte plus cher que l'absence d'une compétence souhaitée.

\section{Expérience}

Le backend ne compte pas simplement la durée totale. Il calcule :

\begin{itemize}
    \item l'expérience totale ;
    \item l'expérience pertinente ;
    \item l'expérience exigée par la fiche de poste.
\end{itemize}

Une expérience de développeur Full Stack ne doit pas automatiquement être considérée comme pertinente pour un poste de Data Analyst.

\section{Responsabilités}

Les responsabilités sont évaluées avec trois niveaux :

\begin{enumerate}
    \item \textbf{preuve complète} : le passage du CV prouve directement la responsabilité ;
    \item \textbf{preuve partielle} : le passage montre une proximité réelle, par exemple reporting ou visualisation de données ;
    \item \textbf{absence de preuve} : aucun passage du CV ne justifie la responsabilité.
\end{enumerate}

Exemple :

\begin{itemize}
    \item \texttt{Power BI} écrit dans le CV donne une preuve forte pour Power BI ;
    \item \texttt{reporting SQL et visualisation de données} donne une preuve partielle pour les tableaux de bord ;
    \item absence de \texttt{SPM} ou \texttt{Foundry} donne 0 sur ces éléments.
\end{itemize}

\section{Redistribution des poids}

Si la fiche de poste ne précise pas la formation ou les langues, ces catégories ne donnent pas automatiquement 100/100. Elles sont ignorées, et les poids sont redistribués sur les catégories réellement présentes.

\chapter{Conclusion}

Le backend SmartRecruit est organisé autour d'une architecture claire :

\begin{itemize}
    \item routes API pour recevoir les fichiers ;
    \item services d'extraction pour transformer les documents en objets structurés ;
    \item normalisation pour rendre les comparaisons cohérentes ;
    \item RAG pour récupérer des preuves dans les CV ;
    \item matching et scoring pour attribuer des scores explicables ;
    \item ranking pour produire un classement final avec gestion des égalités.
\end{itemize}

Cette structure rend le projet évolutif. Le frontend pourra ensuite consommer l'endpoint principal et afficher les résultats sous forme de tableau de classement, avec détails par candidat, compétences manquantes, preuves et explications.

\end{document}
